% Vietnamese translation for OI Public Library
% Source-File: IOI中国国家候选队论文/2003/姜尚仆/姜尚仆.ppt
\documentclass[11pt]{article}
\usepackage{oipl-vn}

\title{Bản trình chiếu: Ứng dụng phương trình tuyến tính modulo}
\author{Jiang Shangpu}
\date{2003}

\begin{document}
\maketitle

\origtitle{Applications of modular linear equations: solving integer problems with number theory}
\sourcepath{IOI China candidate team papers / 2003 / Jiang Shangpu / Jiang Shangpu.ppt}

\section{Phương trình tuyến tính modulo}

Bản trình chiếu bắt đầu từ nhận xét: số học nghiên cứu các tính chất của số
nguyên, nên đặc biệt hữu ích trong những bài có dữ liệu lớn nhưng bản chất
được quyết định bởi chia hết, đồng dư, và chu kỳ.

Dạng phương trình được dùng là
\[
  ax\equiv c\pmod b
\]
hoặc tương đương
\[
  ax+by=c.
\]
Định lý cơ bản trong slide là phương trình có nghiệm khi và chỉ khi
\[
  \gcd(a,b)\mid c.
\]
Khi điều kiện này đúng, thuật toán Euclid mở rộng tìm được một nghiệm riêng,
rồi các nghiệm còn lại được sinh theo modulo thích hợp. Vì thế, sau khi mô
hình hóa được bài toán thành phương trình đồng dư, phần còn lại thường chỉ
là kiểm tra \(\gcd\), \(\operatorname{lcm}\), và số nghiệm trong một miền
hữu hạn.

\section{Bài ví dụ: ball}

Ví dụ chính trong slide là bài \textit{ball}. Một quả bóng xuất phát từ một
ô ở cạnh trái hoặc cạnh dưới của bàn cờ, chuyển động chéo lên. Khi chạm cạnh
bàn, nó phản xạ theo quy tắc gương; khi chạm góc, nó đi ngược lại theo đường
cũ. Câu hỏi là khi quả bóng lần đầu trở lại điểm xuất phát, có bao nhiêu ô
hoặc điểm được đi qua số lần lẻ.

Ta chuyển bài toán hình học thành bài toán trên hình chữ nhật phản chiếu. Gọi
hai kích thước sau khi chuyển đổi là \(L\) và \(H\), và đặt
\[
  g=\gcd(L,H).
\]
Quỹ đạo của bóng là tuần hoàn với chu kỳ đường đi
\[
  2\operatorname{lcm}(L,H).
\]
Vì vậy thay vì mô phỏng từng va chạm, ta xét các khoảng cách ngang và dọc
theo một đường thẳng trong các bản sao của bàn cờ.

\section{Trường hợp chạm góc}

Nếu bóng chạm một góc, nó quay về theo đường cũ và trở lại điểm xuất phát
với hướng ngược lại. Slide viết khoảng cách ngang và dọc dưới dạng
\[
  qL,\qquad pH-a,
\]
nên điều kiện chạm góc trở thành
\[
  qL=pH-a,\qquad\text{hay}\qquad pH-qL=a.
\]
Theo định lý phương trình tuyến tính modulo, điều kiện có nghiệm là
\[
  \gcd(L,H)\mid a.
\]
Trong trường hợp này kết luận của slide là chỉ có hai điểm được đi qua số
lần lẻ.

\section{Trường hợp không chạm góc}

Nếu không chạm góc, bóng trở lại điểm xuất phát với cùng hướng. Các suy luận
trung gian trong slide là:
\begin{itemize}
  \item bóng không thể đi qua một điểm biên quá một lần;
  \item một điểm trong hình chữ nhật được đi qua nhiều nhất hai lần;
  \item số điểm được đi qua số lần lẻ bằng tổng độ dài chu kỳ trừ hai lần số
  điểm được đi qua số lần chẵn.
\end{itemize}

Vấn đề còn lại là đếm các điểm bên trong bị đi qua hai lần. Khi hai lần đi
qua có hướng ngang ngược nhau và hướng dọc giống nhau, giả sử hình chiếu
ngang của điểm là \(x\). Hai khoảng cách ngang có dạng
\[
  2k_1L+x,\qquad 2k_2L-x,
\]
với
\[
  0\le k_1<\frac Hg,\qquad 0<k_2\le \frac Hg.
\]
Điều kiện để chúng tương ứng cùng vị trí dọc là
\[
  (2k_1L+x)-(2k_2L-x)\equiv 0\pmod {2H},
\]
tức
\[
  (k_1-k_2)L+x\equiv 0\pmod H.
\]
Với \(x\) hợp lệ, phương trình theo \(k_1-k_2\) có nghiệm duy nhất modulo
\(H/g\); tổng số cặp thuộc loại này là
\[
  \left(\frac Lg-1\right)\frac Hg.
\]

Tương tự, khi hướng ngang giống nhau và hướng dọc ngược nhau, số cặp là
\[
  \left(\frac Hg-1\right)\frac Lg.
\]

\section{Công thức kết luận}

Kết quả cuối cùng trong slide là:
\[
  \text{nếu } g\mid a,\quad \text{đáp án}=2;
\]
còn nếu \(g\nmid a\), đáp án là
\[
  \frac{2LH}{g}
  -2\left(
    \left(\frac Lg-1\right)\frac Hg
    +\left(\frac Hg-1\right)\frac Lg
  \right).
\]

Điểm quan trọng của lời giải không phải là công thức riêng lẻ, mà là quá
trình đổi mô hình: phản xạ hình học thành chuyển động thẳng, rồi đổi các lần
gặp nhau thành phương trình đồng dư. Phần tổng kết của slide nhấn mạnh ba
điểm khi xử lý bài số học: thiết lập mô hình số học, nắm vững định lý, và
dùng tư duy phân loại để đơn giản hóa bài toán lớn.

\end{document}
